% Zone „Das kennst du schon“ – aus bank/abstaende/zone.jsonl (zusammenbau v0.1)
\einheitenkopf[zone][Kennst du schon]{Das kennst du schon}
\setcounter{aufgabe}{0}

%% TODO Titel: Ich-kann-Satz fehlt in der Bank (Platzhalter: Kettenname) – Nr. 1
%% TODO Anweisung: ein Satz über der ersten Teilaufgabe fehlt in der Bank – Nr. 1
\begin{aufgabe}{Verbindungsvektor bilden (Spitze minus Fuß) und den Betrag als Wurzel aus der Quadratsumme berechnen}
\begin{teile}
\teil A(1 | 2 | 3), B(3 | 3 | 5): $\overrightarrow{AB}$ und seine Länge? $\overrightarrow{AB} = ( \leerfeld | \leerfeld | \leerfeld )$, Länge \leerfeld
\teil C(−1 | 2,4 | 0), D(1 | 0,4 | 1): $\overrightarrow{CD}$ und seine Länge? $\overrightarrow{CD} = ( \leerfeld | \leerfeld | \leerfeld )$, Länge \leerfeld
\end{teile}
\end{aufgabe}

%% TODO Titel: Ich-kann-Satz fehlt in der Bank (Platzhalter: Kettenname) – Nr. 2
%% TODO Anweisung: ein Satz über der ersten Teilaufgabe fehlt in der Bank – Nr. 2
\begin{aufgabe}{Wurzel- und Betragsgleichungen lösen (quadrieren, beide Vorzeichen ansetzen, Lösungen am Sachzusammenhang sieben)}
\begin{gleichungsraster}[2]
\gl{\text{$|x - 3| = 5$ – welche x?}} &
\gl{\text{$|2x - 0{,}5| = 3{,}5$ – welche x?}} \\
\end{gleichungsraster}
\end{aufgabe}

%% TODO Titel: Ich-kann-Satz fehlt in der Bank (Platzhalter: Kettenname) – Nr. 3
%% TODO Anweisung: ein Satz über der ersten Teilaufgabe fehlt in der Bank – Nr. 3
\begin{aufgabe}{Koordinatengleichung einer Ebene lesen und den Normalenvektor ablesen}
\begin{teile}
\teil E: $2x + 3y - z = 4$ – ein Normalenvektor? $\vec n = ( \leerfeld | \leerfeld | \leerfeld )$
\teil E: $0{,}5x + y - 0{,}5z = 2$ – ein ganzzahliger Normalenvektor? $\vec n = ( \leerfeld | \leerfeld | \leerfeld )$
\end{teile}
\end{aufgabe}

%% TODO Titel: Ich-kann-Satz fehlt in der Bank (Platzhalter: Kettenname) – Nr. 4
%% TODO Anweisung: ein Satz über der ersten Teilaufgabe fehlt in der Bank – Nr. 4
\begin{aufgabe}{Geradengleichung in Parameterform lesen und einen Laufpunkt mit Parameter ansetzen}
\begin{teile}
\teil g: $\vec x = (1 | 0 | 2) + t \cdot (2 | 1 | 1)$ – der Punkt für $t = 3$? ( \leerfeld | \leerfeld | \leerfeld )
\teil g: $\vec x = (2 | -1 | 0{,}5) + t \cdot (-1 | 3 | 0)$ – der Punkt für $t = 0{,}5$? ( \leerfeld | \leerfeld | \leerfeld )
\end{teile}
\end{aufgabe}

%% TODO Titel: Ich-kann-Satz fehlt in der Bank (Platzhalter: Kettenname) – Nr. 5
%% TODO Anweisung: ein Satz über der ersten Teilaufgabe fehlt in der Bank – Nr. 5
\begin{aufgabe}{Skalarprodukt berechnen und „null heißt senkrecht“ anwenden}
\begin{teile}
\teil $(1 | 2 | 3) \circ (2 | 0 | 1)$? \leerfeld
\teil $(-0{,}5 | 2 | 4) \circ (2 | 0{,}5 | 1)$? \leerfeld
\end{teile}
\end{aufgabe}

%% TODO Titel: Ich-kann-Satz fehlt in der Bank (Platzhalter: Kettenname) – Nr. 6
%% TODO Anweisung: ein Satz über der ersten Teilaufgabe fehlt in der Bank – Nr. 6
\begin{aufgabe}{Satz des Pythagoras, Satz des Thales und Strahlensätze als Begründungsfiguren}
\begin{teile}
\teil Katheten 6 cm und 8 cm – Hypotenuse? \leerfeld[cm]
\teil Ein 1,6 m großer Mensch wirft 2 m Schatten, ein Baum zur selben Zeit 15 m – wie hoch ist der Baum? \leerfeld[m]
\end{teile}
\end{aufgabe}

%% TODO Titel: Ich-kann-Satz fehlt in der Bank (Platzhalter: Kettenname) – Nr. 7
%% TODO Anweisung: ein Satz über der ersten Teilaufgabe fehlt in der Bank – Nr. 7
\begin{aufgabe}{Verbindungsvektor bilden (Spitze minus Fuß) und den Betrag als Wurzel aus der Quadratsumme berechnen – Fehler finden}
\begin{teile}
\teil Finde den Fehler. A(3 | 1 | 2), B(5 | 7 | 5). \rechnung{\overrightarrow{AB} &= (8 | 8 | 7) \\ |\overrightarrow{AB}| &= \sqrt{64 + 64 + 49} \approx 13{,}30}
\end{teile}
\end{aufgabe}

%% TODO Titel: Ich-kann-Satz fehlt in der Bank (Platzhalter: Kettenname) – Nr. 8
%% TODO Anweisung: ein Satz über der ersten Teilaufgabe fehlt in der Bank – Nr. 8
\begin{aufgabe}{Verbindungsvektor bilden (Spitze minus Fuß) und den Betrag als Wurzel aus der Quadratsumme berechnen – selbst rechnen}
\begin{teile}
\teil C(4 | 2 | 1), D(8 | 6 | 8): $\overrightarrow{CD}$ und seine Länge? $\overrightarrow{CD} = ( \leerfeld | \leerfeld | \leerfeld )$, Länge \leerfeld
\end{teile}
\end{aufgabe}
